%% ma: L11-B15-0062
%% lop: 11
%% chuong: 5
%% bai: 15
%% dang: 11.15.6
%% loai: TN4
%% muc: VD
%% dap_an: "A"
%% nguon: "tai-lieu-hoc-tap-mon-toan-lop-11-theo-chuong-trinh-giao-duc-pho-thong-2018.pdf (Lớp Toán Thầy Duy, Phạm Lê Duy, THPT Chu Văn An), Chương 3 Bài 1, Phần 1, Câu 51 (tr.332)"
%% kien_thuc: "Lớp 11 Bài 15 (giới hạn của dãy số)"
%% trang_thai: cho-duyet
%% ly_do: "Dạng mới đề xuất (11.15.x), chờ giáo viên duyệt dạng; tài liệu gốc không có lời giải mẫu, đáp án do kho tự giải và kiểm bằng sympy."
\begin{ex}
Cho hai dãy số $(u_n)$, $(v_n)$ đều tồn tại giới hạn hữu hạn. Biết rằng hai dãy số đồng thời thỏa mãn các hệ thức $u_{n+1}=4v_n-2$, $v_{n+1}=u_n+1$, $\forall n\in\mathbb Z^+$. Giá trị của giới hạn $\lim(u_n+2v_n)$ bằng
\choice{\True $0$}{$\dfrac32$}{$-1$}{$\dfrac12$}
\loigiai{\sloppy Đặt $\lim u_n=a$, $\lim v_n=b$. Chuyển qua giới hạn: $a=4b-2$, $b=a+1$. Suy ra $a=4(a+1)-2\Rightarrow a=-\dfrac23$, $b=\dfrac13$. Vậy $a+2b=0$.}
\end{ex}
